\chapter{First Outline}
\begin{enumerate}
    \item Introduction
    \item Psychological Networks
    \begin{enumerate}
        \item Network approach to psychopathology
        \item Types of network models
        \begin{enumerate}
            \item Undirected graphs
            \item Directed graphs
        \end{enumerate}
        \item Extending network models
    \end{enumerate}
    \item Moderation
    \begin{enumerate}
        \item Importance of moderation
        \begin{enumerate}
            \item Simpson's paradox
            \item Suppression
        \end{enumerate}
        \item Summary
    \end{enumerate}
    \item Cross-sectional networks
    \begin{enumerate}
        \item Estimation
        \begin{enumerate}
            \item Partial correlation estimation
            \item Nodewise regression
        \end{enumerate}
        \item Example 1: compare groups
        \item Example 2: individual differences
        \item Limitations/future directions
        \begin{enumerate}
            \item Mixed models
            \item Exploratory settings can get complex
            \item Uncertainty about estimation
        \end{enumerate}
    \end{enumerate}
    \item Temporal networks
    \begin{enumerate}
        \item Building blocks of temporal networks
        \item Types of network models from time series
        \item N = 1 networks
        \begin{enumerate}
            \item Seemingly unrelated regression
            \item Generalized least squares
            \item Example
            \item Limitations
        \end{enumerate}
        \item N > 1 networks
        \begin{enumerate}
            \item Considerations--types of networks
            \item GVAR method
            \item Two-stage multilevel method
            \item Example
        \end{enumerate}
        \item Limitations
        \begin{enumerate}
            \item Violating assumptions
            \item Choosing lag length
            \item Interpretation
            \item Alternative approaches: time varying
        \end{enumerate}
    \end{enumerate}
    \item Discussion
    \item Conclusion
\end{enumerate}


%%%%%%%%%%%%%%%%%%%%%%%%%%%%%%%%%%%%
\chapter{Plans for Further Research}
\begin{enumerate}
    \item Discuss some difficulties with network modeling
    \item Causal interpretation
    \item Replicability
    \item Post selection inference
    \item Goal would be to conduct intervention-oriented experiments after observational studies designed as hypothesis generators.
    \item Dynamic structural equation modeling
\end{enumerate}
\subsection{Evaluating the Moderated Network Framework}
\begin{enumerate}
    \item List plans for validating the proposed framework
    \item Specific goals, related to both prediction and inference
    \item Estimating confidence intervals for model parameters.
\end{enumerate}
\subsection{Limitations for Dynamic Networks}
\begin{enumerate}
    \item Lag-order selection
    \item Nonstationarity
    \item Measurement error
    \item Other types of models (e.g., including moving average, cointegration)
    \item Prediction vs. inference
\end{enumerate}
\subsection{Statistical Challenges}
\begin{enumerate}
    \item Separate within-person dynamics from between-person differences
    \item Control for different forms of stability \cite{eichler2012causal} \cite{hamaker2015critique}. Center all level 1 predictors per person \cite{bolger2013intensive,raudenbush2002hlm}. 
    \item Equidistant measurements; correcting when it doesn't hold \cite{boker2010glla}
    \item Lag problem
    \item Dealing with trends \cite{wang2015disaggregating} and cycles in the data \cite{liu2016weekly}.
    \item Multiple time scales \cite{kramer2014time,wichers2014dynamic}.
\end{enumerate}


%%%%%%%%%%%%%%%%%%%%%%%%%%%%%%%%%%%%%%%%
\chapter{Next Outline}
\begin{enumerate}
    \item Introduction
    \begin{enumerate}
        \item Complex systems approach to psychology
        \item Use of network modeling has grown in recent years
        \item However, as these models are still under development, tools don't yet exist for specifying more complex network structures and testing conditional hypotheses
        \item So, in this paper I present moderated network models. This includes some background on what these models are, their statistical foundations, and a description of a new software implementation. 
        \item In this paper I will discuss: (a) what psychological networks are, (b) what moderation models are, (c) how to estimate and interpret moderated networks, (d) how to perform exploratory analyses with model selection procedures, and (e) how these models can be extended further.
    \end{enumerate}
    \item Psychological Networks
    \begin{enumerate}
        \item What are psychological networks?
        \item Types of network models--weighted vs. unweighted; directed vs. undirected
        \item The purpose of network modeling: To understand the conditional dependencies among variables that constitute a psychological system/process.
        \item Limitation of current work is that network models typically only consider pairwise relationships, and not more complex relational structures such as interaction and moderation effects. So, I extend them in this way here.
        \item This current paper will only be dealing with cross-sectional networks, and not longitudinal or time-series networks. Such types of networks will be considered further in the next chapter of my dissertation.
    \end{enumerate}
    \item Multiplicative Interaction Models
    \begin{enumerate}
        \item What is moderation? And what are interaction terms?
        \item Regression framework for testing interactions
        \item The real purpose here is to test conditional hypotheses, including hypotheses about group differences in network structures. 
        \item Criticisms of moderated regression models are often misguided. These models can be powerful when used and interpreted properly. Importantly, we must remember that we are testing conditional effects, and therefore should consider the marginal effects of a given variable with respect to how these might change in accordance with a moderator variable.
        \item Some important considerations are: (a) thinking about the interpretation of the `moderator'. Does theory allow us to rule out the possibility that $M$ has a causal effect on $Y$? Or is the analysis too exploratory to be sure? (b) even if we do identify a moderator, are we interpreting effects that have sufficient support from our sample? That is, what is the observed distribution of the moderator, and how does this relate to the effects that we report? (c) how replicable can we expect our effects to be? Are there other variables that would be important to consider for the context at hand?
        \item In the next section, I will describe how to estimate, analyze, and interpret moderated network models using this framework.
    \end{enumerate}
    \item Estimating Moderated Networks
    \begin{enumerate}
        \item The pcor everyone loves: $pcor(i,j) = -\frac{\kappa_{i,j}}{\sqrt{\kappa_{ii}}\sqrt{\kappa_{jj}}}$. 
        \item And the GGM: $\bSigma = \bDelta\pinv{\I - \bOmega}\bDelta$.
        \item Beginning with undirected networks, we must consider how these models are estimated. The key challenge for psychological network modeling is that we don't typically observe the relations among variables directly. Rather, we estimate the relationships using statistical techniques and weighted graphical representations. 
        \item The two primary ways of estimating network models are: (a) joint estimation, and (b) nodewise regression. The first technique is perfectly fine when we are only considering pairwise relationships. That is, we can estimate the partial correlation matrix of a given dataset by first estimating the covariance matrix, and transforming this into the model of interest. 
        \item With nodewise regression, we estimate the network by testing a series of regression models where each node is the outcome of one model. Then, we aggregate the coefficients across the models in order to estimate the network structure. When there are only main effects in the regression models, we can estimate partial correlations directly from the regression coefficients, even when covariates are included in the models (assuming they are included in all models). When interaction terms are included, however, it is not possible to estimate exact partial correlations, but we can approximate them very closely using terms that have a similar interpretation.
        \item Here's an example of how this would work. Blah blah, estimate each model, then aggregate using the `AND' or `OR' rule to combine coefficients. We can do this for interaction terms as well, although this may not be desirable, and would also require that we model the moderator as a variable in the network. Instead, we may be interested in modeling the moderator as an exogenous variable that influences the network structure but is not affected by it. 
        \item Here's how we can interpret these models. Plots of conditional effects, overall graphical structure, and confidence intervals for estimating marginal effects.
    \end{enumerate}
    \item Model Selection with Moderated Networks
    \begin{enumerate}
        \item Model selection can be a very challenging activity for network modeling, especially when interaction terms are allowed into the model. There are a few ways to approach this, however: (a) Pruning, (b) Criterion-based selection, and (c) Regularization. 
        \item The first one can be done using significance testing and/or false-discovery rate adjustments. The actual models may not necessarily be affected, but the interpretation and visualization of the network can be thresholded in this way. Moreover, the `AND' or `OR' rule can lead to differences in the final estimation of the adjacency matrix, even if the original models stay the same. Another approach could also be to re-estimate the models excluding non-significant terms. This leads us into the second approach to model selection.
        \item Criterion-based selection may mean that we use stepwise testing to find the model that minimizes or maximizes some optimality criterion, such as the AIC, BIC, EBIC, or cross-validated prediction error. Stepwise selection could mean starting with a saturated model and then sequentially dropping non-significant terms until the optimal model is located. This can be extremely time-consuming, however, especially when the number of parameters is large. It also may not be possible to explore the entire model space, and different techniques may lead to different conclusions based on the contour of the probability space.
        \item Regularization is the last and most commonly used approach to model selection in network modeling. The LASSO is the overwhelming favorite within this community, due to its particularly elegant properties. However, it also has some serious drawbacks. One drawback, that also follows with traditional stepwise procedures, is that interaction terms are not taken to have a special kind of status in a model. That is, the model selection procedure can lead to models with higher-order terms and not the constituent lower-order terms. Good reasons to not do this.
        \item Luckily, there is a form of LASSO called the hierarchical LASSO designed specifically to address this problem. We can use this framework for model selection with moderated networks, although it does not necessarily solve all of our problems.
    \end{enumerate}
    \item Model Selection with the Hierarchical LASSO
    \begin{enumerate}
        \item Although the hierarchical LASSO is nice, it does not solve the more general problems of the LASSO. Namely, biased parameter estimates, and the inability to determine the sampling distribution of most parameters. This leaves us with a dilemma.
        \item One response to this problem is to use a two-stage model selection and estimation process. That is, to separate the model selection procedure from the estimation procedure, as is typically done with other approaches. This can allow us to obtain unbiased estimates of the parameters after determining which are most important for the model. 
        \item This doesn't solve all of our problems, though, as the implications of a two-stage modeling process are well-known. The issue is known as `inference after model selection', and is likely one of the most widely known problems in applied statistics. So, more must be done to determine how to deal with this problem.
        \item One possibility is to use this high-dimensional inference framework by one of the LASSO people. This involves a multi-split sampling approach to model selection, along with estimates of confidence intervals that take into account the variation in the estimates as well as their selection frequency. It may not be perfect, but this is likely the best game in town for addressing the problem. 
    \end{enumerate}
    \item Example Analysis
    \item Network Statistics
    \begin{enumerate}
        \item A big issue here is that many people want network-based statistics, such as centrality measures. 
        \item That is fine, but I don't feel super good about it. Anyway, you can still get these with my software.
    \end{enumerate}
    \item P-values and Model Selection
    \begin{enumerate}
        \item FWER Control with single split method; control the family-wise error rate, which is the probability of making at least one false rejection. The data are split into two random, equally sized, disjoint groups, $D_{in} = (X_{in},\,Y_{in})$ and $D_{out} = (X_{out},\,Y_{out})$. Let $\tilde{S}$ be the estimate of the active predictors selected by some variable screening method. Both variable selection and dimensionality reduction are achieved by applying the procedure $\tilde{S}$ to $D_{in}$ (or, the training set). Our goal is to break down the $p$ variables (or edges) down to a smaller number $k \ll p$ with $k$ at most being a fraction of $n$ while keeping all relevant variables. The regression coefficients and $p$-values $\tilde{P}_1,\ldots,\tilde{P}_p$ of the predictors are then determined through $D_{out}$ (the test set) by using OLS given the set of variables $\tilde{S}$ and setting $\tilde{P}_j = 1,\;\forall j \notin \tilde{S}$. If the selected model $\tilde{S}$ contains the true model $S$, i.e., $(\tilde{S} \supseteq S)$, the $p$-values based on $D_{out}$ are unbiased. Finally, each $p$-value $\tilde{P}_j$ is adjusted by a factor $|\tilde{S}|$ to correct for the multiplicity of testing problem. The selected model is given by all variables in $\tilde{S}$ for which the adjusted $p$-value is below a cutoff $\alpha \in (0, 1)$, $\hat{S}_{single} = \big\{ j \in \tilde{S}:\tilde{P}_j |\tilde{S}| \leq \alpha \big\}$. Under suitable assumptions this yields asymptotic control against inclusion of variables in $N$ (false positives) in the sense that: $\limsup\limits_{x\rightarrow\infty}\mathbb{P}\big\lbrack|N\cap\hat{S}_{single}|\geq1\big\rbrack\leq\alpha$.
        \item \textbf{Multi-sample split method}: For each hypothesis, or parameter estimate, a distribution of $p$-values is obtained for random sample splitting. It is proposed that error control can be based on the quantiles of this distribution. For $b = 1,\ldots,B$:
        \begin{enumerate}
            \item Randomly split original data into two disjoint groups $D_{in}^{(b)}$ and $D_{out}^{(b)}$ of equal size.
            \item Using only $D_{in}^{(b)}$, estimate the set of active predictors $\tilde{S}^{(b)}$, which would technically be the $\ell_0$-norm: $||\tilde{S}^{(b)}||_0$.
            \item Using only $D_{out}^{(b)}$, fit the selected variables in $\tilde{S}^{(b)}$ with OLS and calculate the corresponding $p$-values $\tilde{P}_j^{(b)},\;\forall j\in\tilde{S}^{(b)}$. Set the remaining $p$-values to $1$, i.e., $\tilde{P}_j^{(b)} = 1,\ j\notin\tilde{S}^{(b)}$. 
            \item Define the adjusted (non-aggregated) $p$-values as: $\tilde{P}_k^{(b)} = \min\big(\tilde{P}_j^{(b)}|\tilde{S}^{(b)}|,\ 1\big),\ j=1,\ldots,p$.
            \item This leads to a total of $B$ $p$-values for each predictor $j = 1,\ldots,p$. It will turn out that suitable summary statistics are quantiles. For $\gamma\in(0,1)$ define: $Q_j(\gamma) = \min\big\{1,\ q_\gamma\big(\{P_j^{(b)}/\gamma;\ b = 1,\ldots,B\}\big)\big\}$, where $q_\gamma(\cdot)$ is the empirical $\gamma$-quantile function. 
            \item A $p$-value for each predictor $j=1,\ldots,p$ is then given by $Q_j(\gamma)$, for any fixed $0 < \gamma < 1$. This can be shown to be an asymptotically correct $p$-value, adjusted for multiplicity. For instance, with a choice of $\gamma = 0.5$, the quantity $Q_j(0.5)$ is twice the median of all $p$-values $P_j^{(b)},\,b = 1,\ldots,B$. 
            \item Now we have to select for $\gamma$, however. Error control is not guaranteed anymore if we search for the best value of $\gamma$. It is proposed instead to use an adaptive version that selects a suitable value of the quantile based on the data. Let $\gamma_{\min}\in(0,1)$ be a lower bound for $\gamma$, typically $0.05$, and define: $P_j = \min\big\{1,\,(1 - \log\gamma_{\min}) \inf\limits_{\gamma\in(\gamma_{\min},\,1)} Q_j(\gamma)\big\}$. If $\gamma_{\min} = 0.05$, this factor would be upper bounded by $1 - \log(0.05)\approx 3.996$.
        \end{enumerate}
        \item Family-wise error control versus adjustment to false discovery rates. This procedure is family-wise error control using union bound corrections as in Holm. In short, variable $j$ is selected iff $P_j \leq \alpha$, which happens iff $\exists\gamma\in(0.05,1)\ Q_j(\gamma)\leq\alpha/(1-\log0.05)\approx\alpha/3.996$. Equivalently, the $\gamma$-quantile of the adjusted $p$-values, $q_{\gamma}(P_j^{(b)}$, has to be smaller than or equal to $\alpha\gamma/3.996$. This is also the same as the event that the empirical distribution of the adjusted $p$-values $P_j^{(b)}$ for $b=1,\ldots,B$ is crossing above the bound $f(p) = \max\{0.05,\,(3.996/\alpha)p\}$ for some $p\in(0,1)$.
        \item The resulting adjusted $p$-values can be used for both FWER and FDR control. For FWER control at level $\alpha\in(0,1)$, simply all $p$-values below $\alpha$ are rejected and the selected subset is $\hat{S}_{multi} = \{j:P_j \leq \alpha\}$. It can be shown that asymptotically $\mathbb{P}(V > 0) \leq \alpha$, where $V = |\hat{S}_{multi} \cap N|$ is the number of falsely selected variables. Beyond better reproducibility and asymptotic family-wise error control, the multi-split version is more powerful than the single-split method. 
        \item \textbf{FDR control with the multi-split method}. Control of the family-wise error rate is often considered too conservative. If many rejections are made, Benjamini and Hochberg (1995) proposed to control instead for the expected proportion of false rejections, which is the false disovery rate (FDR). Let $V = |\hat{S} \cap N|$ be the number of false rejections for a selection method $\hat{S}$ and $R = |\hat{S}|$ the total number of rejections. The FDR is defined as the expected proportions of false rejections, $\mathbb{E}(Q),\quad\text{with}\ Q = V/\max\{1,R\}$. For no rejections, $R = 0$, the denominator ensures that the false discovery proportion $Q$ is $0$, conforming with the original definition from Benjamini and Bro. 
        \item This approach first orders the observed $p$-values as $P_{(1)} \leq P_{(2)} \leq \ldots \leq P_{(p)}$ and defines $k = \max\{i:P_{(i)} \leq \frac{i}{p}q\}$. Then all variables with the smallest $k$ values are rejected and no rejection is made if the set above is empty. FDR is controlled this way at level $q$ under the condition that all $p$-values are independent. It has been shown that this procedure is conservative under a wide range of dependencies between $p$-values. But maybe this is different with high-dimensional regression? I don't know why it would... In any case, Benjamini and Yekutieli 2001 showed that control is guaranteed at level $q\sum_{i=1}^p i^{-1} \approx q(1/2 + \log p)$. 
        \item The standard FDR is working with the raw $p$-values, for which we assume $P_j \sim \mathcal{U}(0,1)$ for true null hypotheses. The division by $p$ in computing $k$ is an effective correction for multiplicity. The current method already makes an adjustment to the $p$-values, through the penalty of $(1 - \log\gamma_{\min})$, and so further division by $p$ is superfluous. Instead, we can order the corrected $p$-values $P_j,\ j = 1,\ldots,p$ in increasing order $P_{(1)} \leq P_{(2)} \leq \ldots \leq P_{(p)}$ and select the $h$ variables with the smallest p-values, where $h = \max\{i:P_{(i)} \leq iq\}$. The selected set of variables is denoted, with the value of $h$, by $\hat{S}_{multi;FDR} = \{j:P_j \leq P_{(h)}\}$. With no rejections, $\forall i \in \{1,\ldots,p\}\; P_{(i)} > iq \rightarrow \hat{S}_{multi;FDR} = 0$. This procedure will achieve FDR control at level $q\sum_{i=1}^p i^{-1} \approx q(1/2 + \log p)$. To get FDR control at level $q$, we replace $q$ with $q/(\sum_{i=1}^p i\inv)$, completely analogous to the standard FDR procedure under arbitrary dependence of the $p$-values in Benjamini and Hochberg. 
        \item To achieve asymptotic error control, a few assumptions are made in Wasserman and Roeder (2008) regarding crucial assumptions for the variable selection procedure $\tilde{S}$. $A1.\ Screening\ property: \lim_{n \to \infty} \mathbb{P}\lbrack\tilde{S}\supseteq S\rbrack = 1$. $A2.\ Sparsity\ property: |\tilde{S}| < n/2$.
        \item Here are the two error-control procedures:
        \begin{enumerate}
            \item \textbf{FWER}: Assume (A1) and (A2). Let $\alpha,\,\gamma \in (0,1)$. If the null hypothesis $H_{0,j}:\beta_j = 0$ gets rejected whenever $Q_j(\gamma)\leq\alpha$, the family-wise error rate is controlled at level $\alpha$:
            $\limsup\limits_{n\rightarrow\infty}\mathbb{P}\big\lbrack\min\limits_{j\in N} Q_j(\gamma)\leq\alpha\big\rbrack\leq\alpha$. 
            \item \textbf{FDR}: Assume (A1) and (A2). Let $q\in(0,1)$. Let $\hat{S}_{multi;FDR}$ be the set of selected variables and $V = |\hat{S}_{multi;FDR}\cap N|$ and $R = |\hat{S}_{multi;FDR}|$. The FDR with $Q = V/\max\{1,R\}$ is then asymptotically controlled at level $q\sum_{i=1}^p i\inv$, i.e., $\limsup\limits_{n\to\infty}\mathbb{E}(Q)\leq q\sum_{i=1}^p \frac{1}{i}$.
        \end{enumerate}
    \end{enumerate}
    \item Future Directions
    \begin{enumerate}
        \item The next step is to extend this framework to temporal networks with time-series data. This will be the next chapter of my dissertation. A whole new array of issues come up within this framework, and there are many other considerations to take into account.
        \item Again, we need to work on model selection, model comparison, and hypothesis testing.
        \item We also need to think about centrality measures and other network stuff.
        \item Non-normal data and mixed graphical models are other things to consider.
        \item Haslbeck thinks we should look at `moderator centrality', but I'm not so sure. Mostly for the same reason that I'm skeptical of centrality measures more generally.
        \item Information Theory. Shannon entropy denotes the average amount of bits of information (given a set of random variables $\bY$) needed to communicate a discrete outcome. With continuous variables, however, we can use \textit{differential entropy}: $h(\bY) = -\Ev{\log_2 f(\bY)}$, where $f(\bY)$ is the density function of $\bY$. This measure can be computed for any number of variables and quantifies their volatility; for a single variable, this means that the entropy is directly related to the variance. If you then divide $\bY$ into two subsets $\bY^{(1)}$ and $\bY^{(2)}$, we can quantify the association between the two subsets through \textit{mutual information}: $I(\bY^{(1)};\,\bY^{(2)}) = h(\bY^{(1)}) + h(\bY^{(2)}) - h(\bY^{(1)},\bY^{(2)})$. This measure acts as a general measure of the strength of association between any set of variables to any other set of variables. 
        \item Mutual information can be used to look at how closely related two nodes in a network are; how close two sets of nodes in a network are to each other; or even to look at how much mutual information one variable shares with all other nodes in the network (something akin to centrality). When time is involved, we could look at how much information one node has with all other nodes at a subsequent measurement $I(Y_{t1};\,Y_{t+1})$. This could be taken as a temporal centrality measure, which takes into account both the contemporaneous and temporal network. 
        \item When $\bY$ has a multivariate normal distribution with size $P$ and variance-covariance matrix $\bSigma$, as is the case with the GGM and VAR model, the differential entropy becomes: $h(\bY) = \frac{1}{2}\log_2 \big((2\pi e)^P |\bSigma|\big)$. As seen here, this measure is a direct function of the size of the variance-covariance matrix $\bSigma$. This is what allows us to compute all aforementioned mutual information metrics from above. 
        \item For instance, the mutual information between two variables is: $I(Y_i;\,Y_j) = -\frac{1}{2}\log_2 (1 - \sigma_{ij}^2)$. This measure is a direct property of the explained variance $\sigma_{ij}^2$ between two variables. The mutual information of one variable with all other variables is then: $I(Y_i;\,Y_{\setminus i}) = \frac{1}{2}\log_2 \Big(\frac{|\bSigma_{\setminus i}|}{|\bSigma|}\Big)$, where $\bSigma_{\setminus i}$ denotes the variance-covariance matrix without row and column $i$. 
        \item Lastly, in a stationary time series of multivariate normal data, the `temporal closeness' above becomes: $I(y_{t,j};\,\by_{t+1}) = \frac{1}{2}\log_2\Big(\frac{\sigma_{jj}|\bSigma|}{|\bSigma_{j,+}^{TP}|}\Big)$, where $\bSigma_{j,+}^{TP}$ denotes a subsetted Toeplitz matrix: $\bSigma_{j,+}^{TP} = \Var(Y_{ti},\,\bY_{t+1})$.
    \end{enumerate}
    \item Conclusion
    \begin{enumerate}
        \item That's it for now. Again the goal here was to create a new way forward for testing complex hypotheses and fitting moderated network models. 
        \item One contribution is the conceptual idea, and underlying statistical framework. The other contribution, and really the main one, is the software implementation. Much more to do in expanding the generalizing the software, but it currently has many features and is quite flexible. The forthcoming chapters of my dissertation will go into greater detail on this as well. 
    \end{enumerate}
\end{enumerate}

%%%%%%%%%%%%%%%%%%%%%

\chapter{Another Outline}
\begin{enumerate}
    \item Network approach to psychopathology and affective dynamics
    \item Experience sampling approaches and temporal networks
    \item But these approaches haven't taken into account the social situation
    \item Furthermore, there are issues for inference with penalized estimators, which are commonly used in this field.
    \item Therefore, I propose a new method of estimating temporal networks that allows one to specify moderator variables
\end{enumerate}

\section{Network Models in Psychology}
What they are in general. Introduce undirected and directed networks, some background on graph theory, and some applications in psychology. What I'll talk about today are the normal ways that researchers go about fitting network models to psychological data. Good to mention sparsity here as well.
\subsection{Network Approach to Psychopathology}
If we think of the self as a complex system---composed of one's memories, experiences, feelings, and self-perceptions, then we can try and gain a richer picture of how one's predispositions, feelings, and environment interact and direct changes in mood and perhaps even the development of a mood disorder.
\begin{enumerate}
    \item Discuss the general perspective of using networks to understanding psychological disorders
    \item Give examples of some findings
    \item Give examples of personality networks 
    \item Can talk about experience sampling
\end{enumerate}
\subsection{Types of Networks}
\begin{enumerate}
    \item Undirected and Directed models
    \item Distinguish probabilistic graphical models from other kinds of networks
    \item Gaussian Graphical Models and conditional independence
    \item Highlight the utility of partial correlations over associations
    \item Mixed Graphical Models
\end{enumerate}
\subsection{Sparsity Assumption}
\begin{enumerate}
    \item Describe what the sparsity assumption is
    \item Goal is to reduce the number of spurious edges
    \item LASSO is the dominant way this is done in the current literature
    \item It's main function in the network world is to perform variable selection, or rather edge selection. The goal is to reduce the estimation of spurious relationships.
    \item The way it typically works is through determining the ideal lambda parameter value via cross-validation or model selection with the EBIC. 
    \item Implications of the LASSO, though, are that it can improve prediction but makes inference uncertain. Also, it makes specifying moderators difficult
\end{enumerate}

\section{Moderated Networks}
No way to look at moderator variables. Not only has it not been proposed, but the LASSO introduces some major problems. 
\begin{enumerate}
    \item Discuss the importance of moderation in general
    \item Give a clear example
    \item Increases model complexity, but may be crucial for understanding some data.
\end{enumerate}
\subsection{LASSO limitation}
\begin{enumerate}
    \item Describe why the LASSO is problematic
    \item Note that other variable selection methods fall prey to the same problem when the set of possible moderators is not known
\end{enumerate}
\subsection{Other Forms of LASSO}
\begin{enumerate}
    \item Group Lasso
    \item and, most importantly, the hierarchical LASSO
    \item Good to go into some solid detail here
\end{enumerate}

\section{Model Estimation}
Two general approaches: (a) Estimating joint likelihood in terms of inverse covariance matrix, and (b) Performing nodewise regression where each outcome is modeled separately, and the results across models are aggregated to form a network.
\subsection{Undirected Graphs}
These networks are most often used for cross-sectional data. 
\begin{enumerate}
    \item Give some examples
    \item Then talk about how to estimate them
    \item Talk about 'glasso'
    \item Then talk about ML estimation with hard thresholding as an alternative.
    \item Introduce how moderation would work within this context (need to use nodewise regression)
\end{enumerate}
\subsection{Directed Graphs}
\begin{enumerate}
    \item Some examples
    \item Experience sampling data
    \item Unless it's a DAG, it doesn't really have a strong 'network' interpretation. Nevertheless, combining this framework with the time series modeling framework can help represent the results.
    \item Talk about how there are time-varying models being proposed, but these make interpretation somewhat difficult.
\end{enumerate}

\section{Time Series Networks}
\begin{enumerate}
    \item Discuss autoregression
    \item Vector autoregression
    \item Challenges of vector autoregression
    \item Big one is the correlated errors across equations
    \item Also discuss model assumptions
\end{enumerate}
\subsection{SUR Models}
\begin{enumerate}
    \item Seemingly Unrelated Regressions
    \item Generalized least squares
    \item Challenge is performing variable selection and estimating restricted models
\end{enumerate}
\subsection{Graphical VAR}
This is the dominant method being used in the network psychopathology literature.
\begin{enumerate}
    \item Variation on SUR models, using an iterative estimation process that resembles GLS. 
    \item Can get temporal networks and contemporaneous networks
    \item Problem is, can't deal with moderators due to the built-in LASSO to the procedure
\end{enumerate}
\subsection{Two-Step Graphical SUR}
\begin{enumerate}
    \item Perform variable selection before model fitting
    \item Do this in a sequential, nodewise fashion. This has downsides, but also has upsides.
    \item The selected variables for each equation are used to estimate the SUR model with residual covariance matrix.
    \item Discuss advantages over graphical VAR, beyond the accommodation of moderators.
    \item Can be extended to multilevel context---this includes fixed-effect networks and between-subject networks.
\end{enumerate}

\section{Analyzing Network Models}
\begin{enumerate}
    \item Discuss the components that are of interest to researchers
    \item Centrality measures
\end{enumerate}
\subsection{Simulated Example}
Create example from R code. Can also analyze some published data.

\section{Challenges and Future Directions}
\begin{enumerate}
    \item Discuss some difficulties with network modeling
    \item Causal interpretation
    \item Replicability
    \item Post selection inference
    \item Goal would be to conduct intervention-oriented experiments after observational studies designed as hypothesis generators.
    \item Dynamic structural equation modeling
\end{enumerate}
\subsection{Evaluating the Moderated Network Framework}
\begin{enumerate}
    \item List plans for validating the proposed framework
    \item Specific goals, related to both prediction and inference
    \item Estimating confidence intervals for model parameters.
\end{enumerate}
\subsection{Limitations for Dynamic Networks}
\begin{enumerate}
    \item Lag-order selection
    \item Nonstationarity
    \item Measurement error
    \item Other types of models (e.g., including moving average, cointegration)
    \item Prediction vs. inference
\end{enumerate}

\section{Conclusion}
Summarize a bit of what was discussed throughout the paper. Lots still left to do in network modeling, but this offers one potential advance in methodology.